% Generated from validated Article Draft Result. Do not hand-edit; revise the structured draft instead.
\section{Detailed Chronology――一つの『公開日』へ潰さないための2022年史}
\noindent\textbf{年間本文では複数の出来事をstoryとtrajectoryへ圧縮するが、歴史記録ではpaper initial publication、announcement、research preview、beta、general availability、API、open-weight accessを同じ状態として扱えない。本節では、2022年のmaterial eventを後知恵で一本化しないためのchronology規則を明示する。}

研究成果には、論文の初出と、その後の改訂・正式publication・model releaseが別々に存在し得る。年次chronologyではpaper initial publicationを一つの独立状態として扱い、後日のproduct releaseやmodel availabilityの日付で置き換えない。一次資料が日付精度を限定している場合も、推測で時刻や日付を補完しない。


announcement、research preview、closed/open beta、general availability、API提供は、それぞれ利用者が到達できる範囲と責任境界が異なる。DALL-E 2やChatGPTのように一つの名称で複数lifecycle stateを持つ系列は、本文では一つのstoryにまとめても、chronology上では別のeventとして保存する。


open-weightやresearch accessにも段階がある。weightsが取得可能であること、誰でも無条件に取得できること、source codeやtraining recipeが公開されること、license上の再配布・商用利用条件が広いことは同義ではない。OPT、BLOOM、Stable Diffusionなどを比較するときも、各時点のaccess stateを個別に記録する。


このDraft Packageにはchronology専用のevent-level Evidence割当がないため、本節で未検証の日付一覧を新たに生成することはしない。本文各章で検証されたlifecycle境界を、後段のvalidated sourceで独立したeventとして展開するための編集規則をここで固定する。
